% $Header: /home/vedranm/bitbucket/beamer/solutions/conference-talks/conference-ornate-20min.en.tex,v 90e850259b8b 2007/01/28 20:48:30 tantau $

\documentclass{beamer} %[compress]
% alternatively add handout or draft option

%\mode<handout>%{\setbeamercolor{background canvas}{bg=black!5}}
%\usepackage{pgfpages}
%\pgfpagesuselayout{4 on 1}[a4paper,border shrink=5mm,landscape]
% For handout printing, 4 slides on 1 page

\mode<presentation>
{
  \usetheme{Pittsburgh} % plain template Pittsburgh
  \useoutertheme{infolines} % bottom (author, title, place)
  \useoutertheme[subsection=false]{miniframes} % top (sections, frames)
  \usecolortheme{beaver} %blue color whale
	\setbeamertemplate{frametitle}[default][right] % frame titles on the right
  \setbeamercovered{transparent}
  % or whatever (possibly just delete it)
}

\setbeamersize{text margin left=1cm, text margin right=1cm}
%\setbeamertemplate{navigation symbols}{}
\usepackage[english]{babel}
\usepackage{booktabs}
\usepackage{comment}
% or whatever

\usepackage[cp1250]{inputenc}
% or whatever
\usepackage{subfigure} %subfigures

\usepackage{times}
%\usepackage{breqn}
\usepackage{amsmath}
\usepackage[T1]{fontenc}
\def\sym#1{\ifmmode^{#1}\else\(^{#1}\)\fi} %shortcut for Stata tables
% Or whatever. Note that the encoding and the font should match. If T1
% does not look nice, try deleting the line with the fontenc.

%\logo{\includegraphics[height=0.5cm]{logo.jpg}} % logo on every page



\title[Lecture 1: Introduction] % (optional, use only with long paper titles)
{Lecture 1: Introduction to Meta-Analysis}

\author[T.\ Havranek] % (optional, use only with lots of authors)
{ Tomas Havranek
}
% - Give the names in the same order as the appear in the paper.
% - Use the \inst{?} command only if the authors have different
%   affiliation.

\institute[Prague] % (optional, but mostly needed)
{
  	Charles University, CEPR, METRICS

  %\inst%
  %Economic Research Department\\
  %Czech National Bank
      %\inst{1}%
  %Research Department\\
  %Czech National Bank
  %\and
  %\inst{2}%
  %Institute of Economic Studies\\
  %Charles University, Prague
  
%\pgfdeclareimage[height=1cm]{logo}{logo.jpg} % logo only on title page
%\pgfuseimage{logo}
}% - Use the \inst command only if there are several affiliations.
% - Keep it simple, no one is interested in your street address.

\date[Meta in Econ and Finance] % (optional, should be abbreviation of conference name)
{Research Synthesis in Economics and Finance, \\ University of Canterbury\\[2ex]
  February 17, 2025
}%\\[4mm]% - Either use conference name or its abbreviation.
% - Not really informative to the audience, more for people (including
%   yourself) who are reading the slides online


%\subject{Intertemporal Substitution}

\begin{comment}
\AtBeginSection[]
{
  \begin{frame}<beamer>{Outline}
    \tableofcontents[currentsection]%,currentsubsection]
  \end{frame}
}
\end{comment}

\begin{document}

\begin{frame}
  \titlepage
\end{frame}

\begin{comment}

\begin{frame}{Outline}
  \tableofcontents
  % You might wish to add the option [pausesections]
\end{frame}

\end{comment}

%\section{Summary}
%\subsection{}

%\section{EIS}
%\subsection{}


\begin{frame}{Meta important for research, policy, practice}
\textbf{Research experience:}
\begin{itemize}
	\item Professor at Charles University, Prague
	\item Affiliate at the Meta-Research Innovation Center, Stanford
	\item Panel member at ERC (European Research Council) Advanced
\end{itemize}

\bigskip

\textbf{Policy experience:}

\begin{itemize}
	\item Former Advisor to the Board, Czech National Bank
	\item Former member of the Czech National Economic Council
	\item Affiliate at the Centre for Economic Policy Research, London
\end{itemize}

\bigskip

\href{https://tomashavranek.cz/}{\textbf{tomashavranek.cz}}
\end{frame}



\begin{frame}{What happens if the government spends more?}
Labor supply elasticity is key (but how to calibrate it?):
\begin{figure}
	\centering
		\includegraphics[width=1.00\textwidth]{G_shock_Frisch75.pdf}
\end{figure}
\end{frame}



\begin{frame}{Meta in monetary policy making: DSGE models}

DSGE = dynamic stochastic general equilibrium

\bigskip

\begin{enumerate}
	\item The dominant policy analysis and forecasting tool
	\item Used by central banks and ministries of finance
	\item Micro-based with ``deep'' parameters
	\item Immune to Lucas critique (?)
	\item Discipline over intuition
	\item Hard to understand for policy makers
\end{enumerate}
\end{frame}

\begin{frame}{Dozens or hundreds of equations}
\vspace{-0.7cm}
\begin{figure}[htbp]
	\begin{center}
	\subfiguretopcaptrue
	\subfigure{\includegraphics[width=0.495\textwidth]{tonner1}}
	\subfigure{\includegraphics[width=0.45\textwidth]{tonner2}}
		\end{center}
\end{figure}
\end{frame}

\begin{frame}{Psychohistory accomplished? Civ VII?}
\begin{figure}[htbp]
	\begin{center}
\includegraphics[width=1\textwidth]{foundation2.jpg}
		\end{center}
\end{figure}
\end{frame}

\begin{frame}{Main output: fan charts}
\begin{figure}[htbp]
	\begin{center}
\includegraphics[width=1.03\textwidth]{fan2}
		\end{center}
\end{figure}
\end{frame}

\begin{frame}{Calibration uncertainty: meta needed!}
\vspace{-0.4cm}
\begin{figure}[htbp]
	\begin{center}
	\subfiguretopcaptrue
	\subfigure{\includegraphics[width=0.45\textwidth]{dsge1}}
	\subfigure{\includegraphics[width=0.45\textwidth]{dsge2}}
		\end{center}
\end{figure}
\end{frame}

\begin{frame}{What happens if we change one parameter?}
\begin{figure}[htbp]
	\begin{center}
\includegraphics[width=0.95\textwidth]{ir_i}
		\end{center}
\end{figure}
\end{frame}

\begin{frame}{Coincidence? Hardly. Meta-analysis essential}
\begin{figure}[htbp]
	\begin{center}
\includegraphics[width=1\textwidth]{ir}
		\end{center}
\end{figure}
\end{frame}



\begin{frame}{Estimates of the labor supply elasticity vary}
\begin{figure}
	\centering
		\includegraphics[width=0.53\textwidth]{_box2.pdf}
\end{figure}
\end{frame}


\begin{frame}{Values around 0.4--0.5 used for calibration (CBO)}
\begin{figure}
	\centering
		\includegraphics[width=0.9\textwidth]{_hist_frisch.pdf}
\end{figure}
\end{frame}



\begin{frame}{Bias: some estimates less likely to be published}
\begin{figure}
	\centering
		\includegraphics[width=0.9\textwidth]{_caliper2.pdf}
\end{figure}
\end{frame}




\begin{frame}{Context: different situations yield different estimates}
\begin{figure}
	\centering
		\includegraphics[width=0.65\textwidth]{corr_comb2.pdf}
\end{figure}
\end{frame}


\begin{frame}{Implied elasticity: taking bias and context into account}

Let's construct a hypothetical study that uses all information reported in the literature but puts more weight on selected aspects (high precision, quasi-experimental setup, etc.)

\vspace{0.6cm}

\begin{figure}
	\centering
		\includegraphics[width=0.9\textwidth]{best.pdf}
\end{figure}
\end{frame}



\begin{frame}{Course}
\begin{enumerate}
	\item Intro
	\item Applied examples
	\item Literature search
	\item Data collection
	\item Basic meta models
	\item Replications (student presentations)
	\item Publication bias
	\item P-hacking
	\item Heterogeneity
	\item Meta-research
	\item Limitations
	\item Applications (student presentations)
\end{enumerate}
\end{frame}



\begin{frame}{Grading}
\begin{itemize}
	\item Applied meta-analysis (40\%)
	\item Replication (20\%)
	\item Discussion (20\%)
	\item Activity (10\%)
	\item Attendance (10\%)
\end{itemize}
\end{frame}




\begin{frame}{Guidelines}
\begin{figure}
	\centering
		\includegraphics[width=0.9\textwidth]{guide.png}
\end{figure}
\end{frame}



\begin{frame}{Resources}
  \centering
\href{https://meta-analysis.cz/}{\textbf{meta-analysis.cz}}
\end{frame}





\end{document}


